\section{Stand der Forschung}

\subsection*{Agentische Ansätze für Software Engineering}

SWE-bench ist ein Datensatz aus 2\,294 realen GitHub-Issues mit zugehörigen
Codebase-Snapshots, Pull Requests und Test-Suites aus 12 Python-Repositories
\parencite{jimenez2024swebench}.
Der Agent erhält ein Issue und muss einen Patch generieren,
der die zugehörigen fail-to-pass-Tests besteht.
Die Aufgaben erfordern Änderungen in mehreren Dateien, komplexes Reasoning
und Interaktion mit Ausführungsumgebungen.
Klassische Sprachmodelle ohne Werkzeuge schaffen das nicht.

\textcite{yang2024sweagent} zeigen mit SWE-agent, dass ein spezialisiertes
Agent-Computer Interface (ACI, eine strukturierte Schnittstelle zwischen
Modell und Dateisystem) die Lösungsrate auf SWE-bench
von 3,8\,\% (bisherige RAG-Baseline) auf 12,5\,\% anhebt.
SWE-agent steht für den reaktiven Ansatz.
Der Agent entscheidet autonom, welche Dateien er liest,
welche Suchen er ausführt und welche Schritte er unternimmt.

Einen gegensätzlichen Ansatz verfolgt \textcite{xia2024agentless} mit Agentless.
Statt autonomer Exploration verwendet ein dreiphasiger, extern verwalteter Prozess
(Lokalisierung, Patchgenerierung, Testvalidierung) strukturierte Kontextzusammenstellung.
Auf SWE-bench Lite (einer 300-Aufgaben-Teilmenge des vollen SWE-bench)
erreicht Agentless 32\,\% bei nur 0,70\,USD pro Fix.
Das zeigt, dass prozedurale Kontextvorbereitung autonome Agentenexploration schlagen kann.


\subsection*{Wissensgraph-basierte Ansätze}

\textcite{ouyang2024repograph} stellen 2025 RepoGraph vor,
ein Plug-in-Modul zu bestehenden Software-Engineering-Agenten.
Es analysiert den Quellcode via AST (Abstract Syntax Tree,
ein automatisch erstellter Syntaxbaum)
und baut daraus einen Repository-Graphen aus Referenz- und Definitionsknoten
mit Aufruf- und Enthaltensrelationen.
Als Plug-in steigert RepoGraph die Lösungsrate von vier Baselines
(RAG, Agentless, AutoCodeRover, SWE-agent) auf SWE-bench Lite
um durchschnittlich 32,8\,\% relativ gegenüber der jeweiligen Baseline,
ohne die Kernsysteme zu verändern.

\textcite{yang2025kgcompass} entwickeln den graphbasierten Ansatz mit KGCompass weiter. Der Wissensgraph verknüpft nicht nur Code-Struktur,
sondern auch Issues und Pull Requests als Knoten im selben Graphen.
Dadurch lässt sich von einer natürlichsprachigen Bug-Beschreibung
direkt über Mehr-Hop-Pfade zu den relevanten Codestellen traversieren.
KGCompass erreicht 58,3\,\% auf SWE-bench Lite
und eine Fehlerlokalisierungsgenauigkeit von 56,0\,\%.
89,7\,\% der korrekt reparierten Patches erfordern dabei
Mehr-Hop-Verbindungen im Graphen.
Diese Codestellen sind im ursprünglichen Issue nicht direkt erwähnt,
sondern nur über mehrere Abhängigkeitsstufen erreichbar.

\subsection*{Proaktive versus reaktive Kontextbereitstellung}

Die Spannung zwischen proaktiver Vorbereitung und reaktiver Exploration
ist ein eigenständiges Forschungsthema.
\textcite{suri2026codescout} zeigen mit CodeScout,
dass proaktive Kontextanreicherung vor der Agentausführung
der reaktiven Eigenexploration überlegen ist.
Letztere neigt zu Über-Exploration und repetitiven Fix-Mustern,
während CodeScout durch strukturierte Vorab-Exploration präziseren Kontext liefert.

Ergänzend dokumentieren \textcite{li2026contextbench} mit ContextBench
auf 1\,136 Issue-Aufgaben aus 66 Repositories in acht Programmiersprachen:
Aktuelle Coding Agents priorisieren systematisch Recall über Präzision
(es wird mehr abgerufen als nötig),
und es bestehen erhebliche Lücken zwischen dem explorierten
und dem tatsächlich genutzten Kontext.
Komplexere Agent-Architekturen verbessern die Kontextabrufqualität dabei nur marginal.

Das von \textcite{liu2023lostmiddle} dokumentierte Lost-in-the-Middle-Phänomen
gibt diesem Befund eine modellseitige Erklärung.
Selbst als Long-Context-tauglich beworbene Modelle nutzen Informationen
in der Mitte langer Eingaben deutlich schlechter als an Anfang und Ende.
Proaktive, strukturell geordnete Kontextzusammenstellung ist daher
nicht nur eine Effizienzfrage, sondern eine Voraussetzung für zuverlässige Modellnutzung.

\subsection*{Statische Kontext-Dateien als Agenten-Baseline}

Eine weitere Herangehensweise untersucht den Wert manuell gepflegter Kontext-Dateien
(\texttt{CLAUDE.md}, \texttt{AGENTS.md}, \texttt{copilot-instructions.md})
für Coding Agents.
\textcite{gloaguen2026agentsmd} zeigen in einer systematischen Studie auf SWE-bench
(AGENTbench, 138 Aufgaben aus 12 Repositories),
dass vom Entwickler verfasste \texttt{AGENTS.md}-Dateien
die Lösungsrate um durchschnittlich 4\,\% verbessern.
LLM-generierte Dateien dagegen senken sie um 3\,\%.
Beide Varianten erhöhen die Inferenzkosten um mehr als 20\,\%.
\textcite{lulla2026agentsmd} ergänzen auf Basis von 124 Pull Requests aus 10 Repositories,
dass \texttt{AGENTS.md}-Dateien die mediane Bearbeitungszeit um 28,6\,\%
und den Output-Token-Verbrauch um 16,6\,\% reduzieren. 
Das bedeutet, dass Agenten mit besserem Kontext weniger Schritte benötigen, um eine Aufgabe zu lösen.

Diese Befunde belegen: Statische Entwickler-Dokumentation hat einen messbaren,
aber begrenzten Effekt auf Lösungsrate und Effizienz.
Die dargestellten Studien zu Kontext-Dateien vergleichen diese
ausschließlich mit einer Baseline ohne Kontext-Dateien.
Ein direkter Vergleich zwischen statischen Kontext-Dateien
und graphbasierter Kontextzusammenstellung findet sich
in der hier zitierten Literatur nicht.

\subsection*{Transparenz von Coding Agents}

Eine weitere offene Frage betrifft die Inspizierbarkeit agentischer Prozesse.
\textcite{agentstepper2026} stellen AgentStepper vor,
einen interaktiven Debugger für Coding Agents (SWE-agent, ExecutionAgent, RepairAgent),
der Entwicklern ermöglicht, Ausführungs-Trajektorien
(die Abfolge aller Schritte, die ein Agent bei einer Aufgabe durchführt)
schrittweise zu durchsuchen, Prompt-Verläufe einzusehen
und Agenten-Entscheidungen nachzuverfolgen.
In einer kontrollierten Studie mit 12 Teilnehmenden
stieg die Bug-Erkennungsrate von 17\,\% auf 60\,\%.
Diese Arbeit belegt, dass Einblick in die Ausführungsschritte
die Fähigkeit von Entwicklern zur Fehlerkorrektur messbar verbessert.

AgentStepper zeigt, welche Tool-Calls der Agent ausgeführt hat
und welche Ergebnisse sie lieferten.
Das im Rahmen dieser Arbeit zu entwickelnde System würde damit
auch über AgentStepper inspizierbar sein,
da seine Ergebnisse explizite Begründungen der Kontextauswahl enthalten.
Warum wurde genau diese Codestelle inkludiert?
Über welche Graph-Kante wurde sie erreicht?
Warum wurde eine andere Stelle ausgeschlossen?
KGCompass bettet Entscheidungspfade als Prompt-Kontext ein,
um die Patch-Generierung des LLM zu verbessern \parencite{yang2025kgcompass}.
Die Pfade dienen der Modellqualität, nicht der Entwicklerkontrolle.
Ob ein Entwickler die Kontextauswahl vor der Agentausführung bewerten
und korrigieren kann, wird nicht untersucht.


\subsection*{Integrationsinfrastruktur: MCP}

Im November 2024 hat Anthropic das Model Context Protocol (MCP)
als offenen Standard veröffentlicht \parencite{anthropic2024mcp}.
MCP definiert ein JSON-RPC-basiertes Protokoll zwischen einem Client (z.\,B. Claude Code)
und einem Server, der Kontext bereitstellt.
Ein Server kann drei Typen von Primitiven anbieten:
ausführbare Funktionen (Tools),
Datenquellen wie Dateiinhalte oder Datenbankrecords (Resources)
und wiederverwendbare Prompt-Templates auf dem Server (Prompts).
Für diese Arbeit ist MCP die natürliche Integrationsschicht:
Der Living-Knowledge-Graph-Server stellt seinen \texttt{get\_context}-Aufruf
als MCP-Tool bereit und ist damit ohne Änderung am Agent
in Claude Code, Cursor und kompatible Clients einbindbar.

\subsection*{Was diese Arbeit untersucht}

Die dargestellten Arbeiten belegen:
Graphbasierte Kontextzusammenstellung verbessert die Agent-Performance messbar.
Proaktive Kontextbereitstellung schlägt reaktive Exploration in Präzision und Effizienz.
Statische Entwickler-Dokumentation hat einen messbaren, aber begrenzten Effekt.
Diese Befunde basieren ausnahmslos auf Benchmark-Evaluationen auf SWE-bench.
SWE-bench-Aufgaben basieren auf eingefrorenen Codebase-Snapshots mit vorgefertigten Testsuiten.
Der Mehrwert eines kontinuierlich synchronisierten Wissensgraphen lässt sich dort nicht messen.
Laufende Hochschulprojekte werden aktiv weiterentwickelt, haben keine vollständige Testabdeckung
und involvieren menschliche Entwickler mit unterschiedlichem Erfahrungsstand.
Ob ein Living Knowledge Graph auch unter diesen Bedingungen einen messbaren Mehrwert liefert,
ist bislang nicht untersucht.

Daraus ergeben sich drei Fragen:

Wie lässt sich ein Wissensgraph dauerhaft mit einer aktiv weiterentwickelten
Codebase synchron halten?
Bestehende Systeme analysieren das Repository zum Aufgabenzeitpunkt;
eine inkrementelle, commit-getriebene Aktualisierung fehlt.

Verbessert ein Living Knowledge Graph die Arbeitseffizienz von Studierenden
an realen Softwareprojekten messbar,
verglichen mit reaktiver Exploration und statischer Dokumentation?

Ermöglicht die explizite Graph-Pfad-Begründung im MCP-Response Entwicklern,
Kontextentscheidungen des Agenten vor der Ausführung zu bewerten
und bei Bedarf zu korrigieren?

Die vorliegende Arbeit beantwortet diese Fragen mit einem Living Knowledge Graph,
der inkrementell mit der Codebase synchronisiert wird
und als MCP-Server in Claude Code, Cursor und kompatible Umgebungen einbindbar ist.
Jeder gelieferte Kontext enthält einen expliziten Entscheidungspfad,
sodass Entwickler die Kontextauswahl nachvollziehen und bei Bedarf eingreifen können.
